\section{The integration layer}
\label{sec:method}

The suite does not fork the microsimulation; it drives the released \texttt{policyengine} package through a thin adapter module (\texttt{policyengine\_macro.core}) whose job is to validate inputs, run the model, and return a payload with declared provenance and a common cross-model score block. Because importing the package pulls in the full UK and US country models---roughly twenty seconds---the import is deliberately lazy, inside the adapters rather than at module level, so that argument validation fails fast even where the model is not installed.

\subsection{Household adapters}

\texttt{pe\_household} evaluates one household under a single policy environment and returns the resolved person-, benefit-unit-/tax-unit- and household-level variables together with a headline summary; \texttt{pe\_household\_impact} runs the same household twice, with and without the reform, and returns baseline, reform, and their differences---the $\Delta(h)$ of equation~\eqref{eq:hhdelta}. Neither touches any dataset, and neither requires a credential.

Two input guards are worth documenting because both were written in response to failure modes that produce a plausible wrong answer rather than an error. First, \texttt{country} is required and has no default on the household tools: the UK and US are different tax--benefit models, and a default would silently score the wrong country for a US household. (\texttt{pe\_population\_impact} retains a \texttt{country="uk"} default for backwards compatibility---an asymmetry the code declares in a comment and Section~\ref{sec:limitations} restates as a limitation rather than a feature.) Second, the US household dictionary is validated against the key the model actually reads: PolicyEngine US takes the state as household-level \texttt{state\_code\_str} and falls back to California when it is absent, so a misspelled \texttt{state\_code} previously returned a confident California answer for a Texas household. The adapter now refuses the wrong key by name, and refuses a state code that is not a two-letter US state or DC.

\subsection{The population pipeline}
\label{sec:poppipeline}

Population scoring runs the same rule graph over representative household microdata. The pipeline is: obtain the dataset for the requested country, year and vintage (downloading or building it on first use); construct and run a baseline simulation; construct and run a counterfactual differing only in the policy; and measure the difference. The baseline simulation is cached in-process per \texttt{(country, year, dataset)}, so a session that scores several reforms against the same baseline pays for one baseline and one reform run each thereafter.

The measurement core is deliberately shared between two callers that mean different things. \texttt{pe\_population\_impact} scores a \emph{reform}: policy differs between the two runs. \texttt{\_pe\_population\_incidence} scores a \emph{macro shock}: policy is identical on both sides and the counterfactual differs only by an overlay scaling the employment-income inputs, so the budgetary difference it reports is the automatic-stabiliser response to the shock and not the cost of anything. Both report the same outcome block, and the second is careful to say in its own headline that policy was unchanged. When the overlay is absent the shocked run is construction-identical to the baseline and every delta is zero by design, which callers surface as an honest null rather than an error.

Three measurement choices in that shared core are stated here because they determine what the published aggregates mean.

\paragraph{The budget measure differs by country.} For the UK the budgetary impact is the change in \texttt{gov\_balance}---all modelled taxes, including capital gains tax and employer National Insurance, minus benefit spending---because that is the UK budget headline and it is not in the model's default output set, so the adapter requests it explicitly. For the US it is the change in \texttt{household\_tax} minus the change in \texttt{household\_benefits}. These are different concepts on different coverage, and the returned payload names the basis it used in a \texttt{budgetary\_impact\_basis} field rather than leaving the reader to assume they are comparable.

\paragraph{Decile effects are means, not means of ratios.} Distributional output groups households by baseline income decile and reports the change in each decile's mean net income, with the relative figure computed as the change in the decile's mean expressed as a percentage of that mean. The obvious alternative---averaging each household's own percentage change---is dominated by households with near-zero baseline income, for whom a small absolute change is an enormous ratio, and is not used.

\paragraph{Winner and loser counts equal the sum of the rows.} The headline counts accumulate exactly the integers displayed in the decile rows. Truncating the underlying floats independently at each level had let the headline disagree with the table beneath it.

\subsection{Provenance and the common score block}

Every adapter returns a provenance record naming the model identifier, the installed distribution and its version, the source repository, the data vintage, and the baseline definition, together with an instruction for reproducing the run---check out the recorded source at the recorded revision, install the recorded adapter version, and rerun the serialised inputs. For population runs the data vintage \emph{is} the dataset name, which is how the registry's ``recorded per run'' vintage policy is implemented (Section~\ref{sec:repro}).

Population scores additionally carry a \texttt{ScoreResult} block in the schema shared across the suite, so that the same reform scored through the microsimulation, the OBR emulator and the OLG model can be compared field by field. Three of its fields are load-bearing for this paper's argument. \texttt{assumptions} always contains ``static microsimulation: no behavioural or macro feedback'' and the dataset name with its household count. \texttt{caveats} always contains a statement that GDP, consumption and investment are out of scope for a static microsimulation and only the budgetary impact is filled---so a consumer reading the common schema cannot mistake an empty GDP field for a zero effect. And the revenue quantity's \texttt{uncertainty} field reads, verbatim, ``not estimated in this result.'' That is the machine-readable form of Section~\ref{sec:uncertainty}'s central admission.

\subsection{Hosted surfaces}

The model is reachable three ways: a hosted Model Context Protocol server exposing \texttt{calculate\_household}, \texttt{household\_reform\_impact}, \texttt{population\_reform\_impact}, \texttt{list\_reform\_parameters} and \texttt{score\_reform} with \texttt{model="microsim"}; the \texttt{pe-macro} command-line interface mirroring the same tools; and the Python package directly. Runtime is sub-second for a household and minutes for a population run. On the hosted server the gated UK microdata credential is provisioned server-side, so a hosted population score requires nothing of the caller---which is convenient, and is also the reason Section~\ref{sec:repro} distinguishes between a result being \emph{obtainable} and a result being \emph{reproducible by the reader}.
